% Part: first-order-logic
% Chapter: natural-deduction
% Section: derivations

\documentclass[../../../include/open-logic-section]{subfiles}

\begin{document}

\iftag{FOL}
      {\olfileid{fol}{ntd}{der}}
      {\olfileid{pl}{ntd}{der}}

\olsection{推导}

\begin{explain}
我们已经说明了假设是什么，并给出了推理规则。自然演绎中的推导是由它们归纳生成的：每个推导要么本身是一个假设，要么由一个、两个或三个推导后跟一个正确的推理构成。
\end{explain}

\begin{defn}[推导]
\Article{derivation} 从假设~$\Gamma$ 出发对语句~$!A$ 的一个\emph{推导}，是满足下列条件的有限语句树：
\begin{enumerate}
\item 树中最顶层的语句要么属于~$\Gamma$，要么被树中的某个推理所解除。
\item 树中最底层的语句是~$!A$。
\item 树中除最底层的语句~$!A$ 外，每个语句都是某次推理规则正确应用的前提，且该次应用的结论在该树中紧接于该语句下方。
\end{enumerate}
此时，我们称~$!A$ 为该推导的\emph{结论}，称~$\Gamma$ 为其未解除的假设。

如果存在从~$\Gamma$ 到~$!A$ 的推导，我们就说~$!A$ \emph{可由~$\Gamma$ 推导}，记作 $\Gamma \Proves !A$。如果存在~$!A$ 的一个推导，其中所有假设都被解除了，我们就记作~$\Proves !A$。
\end{defn}

\begin{ex}
每个假设本身就是一个推导。因此，例如，$!A$ 本身是一个推导，$!B$ 本身也是一个推导。我们可以通过应用（比如说）$\Intro{\land}$ 规则，从这些推导得到一个新的推导：
\begin{prooftree}
\AxiomC{$!A$}
\AxiomC{$!B$}
\RightLabel{\Intro{\land}}
\BinaryInfC{$!A \land !B$}
\end{prooftree}
这些规则是通用的：我们可以用任何语句替换其中的 $!A$ 和~$!B$，例如换成 $!C$ 和~$!D$。那么结论就会是 $!C \land !D$，因此
\begin{prooftree}
\AxiomC{$!C$}
\AxiomC{$!D$}
\RightLabel{\Intro{\land}}
\BinaryInfC{$!C \land !D$}
\end{prooftree}
是一个正确的推导。当然，我们也可以交换假设，让 $!D$ 扮演~$!A$ 的角色，让 $!C$ 扮演~$!B$ 的角色。因此，
\begin{prooftree}
\AxiomC{$!D$}
\AxiomC{$!C$}
\RightLabel{\Intro{\land}}
\BinaryInfC{$!D \land !C$}
\end{prooftree}
也是一个正确的推导。

我们现在可以应用另一条规则，比如 $\Intro{\lif}$。该规则允许我们得出一个条件句作为结论，并允许我们解除任何与该条件句前件相同的假设。所以以下两个都是正确的推导：
\begin{prooftree}
\AxiomC{$\Discharge{!C}{1}$}
\AxiomC{$!D$}
\RightLabel{\Intro{\land}}
\BinaryInfC{$!C \land !D$}
\DischargeRule{\Intro{\lif}}{1}
\UnaryInfC{$!C \lif (!C \land !D)$}
\DisplayProof\bottomAlignProof
\AxiomC{$!C$}
\AxiomC{$\Discharge{!D}{1}$}
\RightLabel{\Intro{\land}}
\BinaryInfC{$!C \land !D$}
\DischargeRule{\Intro{\lif}}{1}
\UnaryInfC{$!D \lif (!C \land !D)$}
\end{prooftree}
它们分别证明了 $!D \Proves !C \lif (!C \land !D)$ 和 $!C \Proves !D \lif (!C \land !D)$。

请记住，解除假设是允许的，而非必须的：我们并非必须解除假设。
特别地，即使假设并未出现在推导中，我们也可以应用相应的规则。例如，下面的推导是合法的，尽管并不存在需要被解除的假设~$!A$：
\begin{prooftree}
  \AxiomC{$!B$}
  \DischargeRule{\Intro{\lif}}{1}
  \UnaryInfC{$!A \lif !B$}
\end{prooftree}
\end{ex}

\end{document}